\providecommand{\StudyOptions}{}
\expandafter\documentclass\expandafter[\StudyOptions]{lumenbook}
\title{A structured PDE notebook}
\author{Lumenbook}
\begin{document}
\chapter{Study with structure}
\readingsource{Original template example / 模板原创示例}
  {Heat equation: one model, several ways to understand it.}
\begin{learningobjectives}
  \item Identify the equation, initial data and boundary values.
  \item Relate a sine mode to an energy estimate.
  \item Reconstruct a proof before opening the solution.
\end{learningobjectives}
\marginprogress{3}{8}
\section{References that know their type}
The \term{energy method}{能量方法} turns a differential equation into
an estimate for an integral quantity.
\begin{definition}[Energy]\label{def:study-energy}
For a sufficiently smooth function $u$, set
$E(t)=\frac12\int_0^L |u(x,t)|^2\dd x$.
\end{definition}
\begin{theorem}[Dissipation]\label{thm:study-energy}
For a classical heat solution with zero endpoint values,
$E'(t)=-\int_0^L|u_x(x,t)|^2\dd x\leq 0$.
\end{theorem}
\begin{lemma}[A sine mode]\label{lem:study-mode}
If $k=\pi/L$ and $u(x,t)=a(t)\sin(kx)$, then $u_t=u_{xx}$
is equivalent to $a'(t)=-k^2a(t)$.
\end{lemma}
\Cref{def:study-energy} introduces the quantity used in
\cref{thm:study-energy}. The reduction in \cref{lem:study-mode}
suggests an explicit example. Labels keep their types while sharing
one chapter-based numerical sequence.

\newpage
\section{Practice before looking at the answer}
\begin{exercise}[The first mode]\label{ex:study-mode}
Solve $u_t=u_{xx}$ on $0<x<L$ with zero endpoint values and
$u(x,0)=\sin(\pi x/L)$. Compute the decay rate of $E(t)$.
\end{exercise}
\begin{hint}[ex:study-mode]
Start with the separated form $u(x,t)=a(t)\sin(\pi x/L)$.
Apply \cref{lem:study-mode} and use $a(0)=1$.
\end{hint}
\begin{solution}[ex:study-mode]
The time factor decays exponentially. Integrating
$a'(t)=-k^2a(t)$ with $a(0)=1$ gives
\begin{equation}\label{eq:study-answer}
  u(x,t)=e^{-k^2t}\sin(kx),\qquad k=\frac{\pi}{L}.
\end{equation}
The energy is $E(t)=\frac{L}{4}e^{-2k^2t}$.
\begin{codeblock}[language=Python,title={Decay rate}]
from math import pi
decay_rate = (pi / L)**2
energy_rate = 2 * decay_rate
\end{codeblock}
\end{solution}
\begin{pitfall}
The amplitude and energy decay at different rates. Squaring the
amplitude doubles the exponent. 注意区分振幅与能量的衰减率。
\end{pitfall}

\newpage
\section{A deliberate review}
\begin{chaptersummary}
  \item A boundary condition can remove an integration-by-parts term.
  \item A separated solution provides a concrete check of an estimate.
  \item Energy decay describes a global quantity, not pointwise monotonicity.
\end{chaptersummary}
\begin{reviewchecklist}
  \checkitem[done]{I can state the equation and all data precisely.}
  \checkitem[done]{I can recover the sine-mode calculation without notes.}
  \checkitem[question]{I still need to justify the regularity assumptions.}
  \checkitem{I can adapt the argument to a new boundary condition.}
\end{reviewchecklist}
\begin{readingbox}{Progress is supplied by the reader}
Use the bar to track a chosen page range, exercise list or reading plan.
The values are entered manually; completing pages is distinct from
understanding their proofs.
\readingprogress{3}{8}
\end{readingbox}
\begin{codeblock}[language=LaTeX,title={Two editions from one source}]
% Reading edition: hints and solutions are visible by default.
\documentclass{lumenbook}
% Self-test edition:
% \documentclass[solutions=false,hints=false]{lumenbook}
\end{codeblock}
\end{document}
